\documentclass[11pt]{amsart}
\usepackage[margin=1.1in]{geometry}
\usepackage{amsmath,amssymb,amsthm}
\usepackage{booktabs}
\usepackage[colorlinks=true,linkcolor=blue,citecolor=blue,urlcolor=blue]{hyperref}

\newtheorem{theorem}{Theorem}
\newtheorem{lemma}[theorem]{Lemma}
\newtheorem{corollary}[theorem]{Corollary}
\theoremstyle{remark}
\newtheorem{remark}[theorem]{Remark}

\newcommand{\lean}[1]{\texttt{#1}}
\setlength{\emergencystretch}{3em}

\begin{document}

\title[Two conjectures of Sun on A217703]{Two conjectures of Sun on the recurrence $a(n+1)=2n(n+1)a(n)-n^4a(n-1)$ are false}

\author{Ansh Mishra}
\author{Aryan Senthilkumar}

\date{October 2026}

\begin{abstract}
Let $a(0)=1$, $a(1)=0$ and $a(n+1)=2n(n+1)a(n)-n^4a(n-1)$ (OEIS A217703).
Z.-W. Sun conjectured that (ii) $a(n)<0$ exactly when $1<n<58$ or $n>2177$, and that (iii) $|a(n)|^{1/n}=o(n^2)$.
Both are false.
An exact computation shows $a(81935)>0$.
We also prove that $\limsup_n|a(n)|^{1/n}/n^2=e^{-2}$, using the Casoratian of the recurrence and a polynomial bound for a second solution.
The disproof of (iii) is formalized in Lean~4 with Mathlib, against the statement \lean{OeisA217703.conjecture3} in Google DeepMind's Formal Conjectures repository.
\end{abstract}

\maketitle

\section{Introduction}

The sequence $a(n)=1,0,-1,-12,-207,\dots$ is A217703 \cite{OEIS}.
It is the constant term $S_n(0)$ of polynomials $S_n(x)$ that Sun introduced in 2013.
In the OEIS entry Sun conjectured, among other things:
\begin{itemize}
\item[(ii)] $a(n)<0$ if and only if $1<n<58$ or $n>2177$;
\item[(iii)] $|a(n)|^{1/n}=o(n^2)$ as $n\to\infty$.
\end{itemize}
They are stated in \cite{FC} as \lean{OeisA217703.conjecture2} and \lean{conjecture3}.

\begin{theorem}\label{thm}
\begin{enumerate}
\item[(a)] $a(81935)>0$; hence (ii) is false.
\item[(b)] $\limsup_{n\to\infty}|a(n)|^{1/n}/n^2=e^{-2}$; hence (iii) is false.
\end{enumerate}
\end{theorem}

\section{Proof of (a)}

We evaluated the recurrence exactly with arbitrary-precision integers.
The signs of $a(n)$ change at $n=2,58,2178,81935$; the next change, located with $40$-digit floating point, is at $3082582$.
In particular $a(81935)$, an integer of about $2.43$ million bits, is positive.
The successive ratios $2178\to81935\to3082582$ are close to $e^{2\pi/\sqrt3}\approx37.6$.
This matches the indicial equation $\alpha^2+\alpha+1=0$ of the normalized recurrence below, whose solutions behave like $n^{-1/2}\cos(\frac{\sqrt3}2\log n+\phi)$.
This heuristic only explains the pattern; the disproof is the exact computation.

\section{Proof of (b)}

Put $b(n)=a(n)/(n!)^2$.
For $n\ge1$,
\begin{equation}\label{eq:b}
(n+1)^2b(n+1)=2n(n+1)\,b(n)-n^2\,b(n-1),\qquad b(0)=1,\ b(1)=0 .
\end{equation}
Let $v$ be the solution of \eqref{eq:b} with $v(0)=0$, $v(1)=1$.

\begin{lemma}\label{lem:poly}
$|v(n)|\le n$ and $|b(n)|\le n$ for all $n\ge1$.
\end{lemma}

\begin{proof}
Let $u$ be either solution, and put $s(n)=n\,u(n)$ and $\delta(n)=s(n)-s(n-1)$.
For $n\ge2$, \eqref{eq:b} becomes $s(n+1)=2s(n)-s(n-1)-\frac{s(n-1)}{n^2-1}$, that is,
\[
\delta(n+1)=\delta(n)-\frac{s(n-1)}{n^2-1}.
\]
We claim that $|s(k)|\le k^2$ for $1\le k\le n$ and $|\delta(k)|\le k$ for $2\le k\le n$.
This holds for $n=2$, because $(s(1),s(2))$ is $(1,2)$ for $v$ and $(0,-\frac12)$ for $b$.
Inductively,
\[
|\delta(n+1)|\le n+\frac{(n-1)^2}{n^2-1}<n+1\qquad\text{and}\qquad|s(n+1)|\le n^2+n+1\le(n+1)^2 .
\]
Hence $|u(n)|=|s(n)|/n\le n$.
\end{proof}

\begin{lemma}[Casoratian]\label{lem:cas}
$b(n)v(n-1)-b(n-1)v(n)=-1/n^2$ for all $n\ge1$.
\end{lemma}

\begin{proof}
The value at $n=1$ is $-1$.
Multiply \eqref{eq:b} for $b$ by $v(n)$, multiply \eqref{eq:b} for $v$ by $b(n)$, and subtract.
This gives $(n+1)^2W(n+1)=n^2W(n)$, where $W(n)=b(n)v(n-1)-b(n-1)v(n)$.
\end{proof}

\begin{proof}[Proof of (b)]
By the two lemmas, for $n\ge2$,
\[
\frac1{n^2}\le|b(n)|(n-1)+|b(n-1)|\,n\le(2n-1)\max\bigl(|b(n)|,|b(n-1)|\bigr).
\]
So for each $n\ge2$ some $m_n\in\{n-1,n\}$ has $|b(m_n)|\ge1/(2n^3)$.
The set $M$ of these $m_n$ is infinite.
For $m\in M$, using $m!\ge(m/e)^m$,
\[
\frac{|a(m)|^{1/m}}{m^2}=\frac{(m!)^{2/m}\,|b(m)|^{1/m}}{m^2}\ \ge\ e^{-2}\bigl(2(m+1)^3\bigr)^{-1/m}\ \longrightarrow\ e^{-2}.
\]
Conversely, $|b(n)|\le n$ and $(n!)^{2/n}\le(n+1)^2/e^{2}\cdot(e\,(n+1))^{2/n}$ give $\limsup\le e^{-2}$.
\end{proof}

\section{Formalization and checks}

The negation of \lean{OeisA217703.conjecture3} is formalized in Lean~4 with Mathlib as \lean{OeisA217703.conjecture3\_false}.
Its type is literally $\neg$ (the repository statement), which is checked by \lean{rfl}.
The formal proof works with integers.
It uses the Casoratian identity $a(k+1)w(k)-a(k)w(k+1)=-(k!)^4$ for the second solution $w=(n!)^2v$, a crude bound $|w(n)|\le3^n(n!)^2$, and $n^n\le3^nn!$.
From these it shows that $|a(n)|^{1/n}/n^2\ge1/108$ for infinitely many $n$.
The proof depends only on the axioms \lean{propext}, \lean{Classical.choice} and \lean{Quot.sound}.
Part~(a) is a pure computation with numbers of millions of bits, and we have not formalized it.
Our scripts check Lemmas~\ref{lem:poly} and \ref{lem:cas} with exact rationals for $n<3000$ and compute the signs of $a(n)$ exactly for $n\le81940$.

\begin{thebibliography}{9}

\bibitem{FC}
Google DeepMind, \emph{Formal Conjectures}, \url{https://github.com/google-deepmind/formal-conjectures}; the statements are in \texttt{FormalConjectures/OEIS/217703.lean}, pull request 6426 (2026).

\bibitem{OEIS}
OEIS Foundation Inc., \emph{The On-Line Encyclopedia of Integer Sequences}, \url{https://oeis.org/A217703}.

\end{thebibliography}

\end{document}
